\documentclass[10pt,a4paper]{amsart}
\usepackage[margin=19mm]{geometry}
\usepackage{amsmath,amssymb,mathtools}
\usepackage[scr=rsfs,cal=euler]{mathalfa}
\usepackage{xfrac}
\usepackage[unicode,hidelinks,bookmarks=false]{hyperref}
\newcommand{\ie}{i.e.\ }
\newcommand{\cf}{cf.\ }
\newcommand{\resp}{respectively\ }
\newcommand{\Subsection}{\subsection}
\providecommand{\nobreakdash}{\nobreak\hbox{-}}
\setlength{\emergencystretch}{2em}
\multlinegap0pt
\usepackage{bm}
\usepackage{bbm}
\usepackage{dsfont}
\usepackage{xspace}
\usepackage{cancel}
\usepackage{mathtools}
\usepackage{amsthm}
\makeatletter
\@ifundefined{Thm}{\newtheorem{Thm}{Theorem}}{}
\makeatother
\title[Singularity for graph-directed conjugate equations indexed b]{Singularity for graph-directed conjugate equations indexed by a two-vertex digraph}
\author{Kazuki Okamura}
\date{Source: arXiv:2603.21536v1 (CC BY 4.0)}
\begin{document}
\maketitle

\noindent\emph{Source and licence.} Singularity for graph-directed conjugate equations indexed by a two-vertex digraph. Version arXiv:2603.21536v1 (CC BY 4.0), distributed under the Creative Commons Attribution 4.0 International licence, as recorded in the arXiv licence metadata for this exact version and shown on the abstract page https://arxiv.org/abs/2603.21536v1. Full rights notice: licenses/arxiv-2603.21536-CC-BY-4.0.txt.\par\smallskip

\noindent\textit{Editorial scope.} Counted unit: the Thm whose source label is \texttt{thm:LF} in the article (source0.tex lines 357--366, source label \texttt{thm:LF}), delivered together with the article's own complete original proof of it (source0.tex lines 368--496, one \texttt{proof} environment of 129 lines). No mathematics is written, repaired or re-derived: statement and proof are the article's own text line for line. Required context retained uncounted: eq:gd-dR-fe-def-alt (lines 178--180); eq:def-f-dyadic (lines 326--328); eq:cd-posi (lines 334--336); eq:compatible-LF (lines 338--340); eq:def-alpha (lines 344--346). Every definition, display and label of those lines is retained unchanged. Omitted from this package and not redistributed: 1--177, 181--325, 329--333, 337--337, 341--343, 347--356, 367--367, 497--675. Nothing inside a retained excerpt is omitted or altered: every retained body is delivered exactly as the article's lines are written, so no omission receipt of any kind accompanies this package. Citations: the retained excerpts carry no citation command at all, so no bibliography entry of the article is transcribed and no reference is redistributed. Reference adaptation: none was needed and none was made; every reference inside the delivered unit resolves inside it, so no unresolved reference or citation is produced. Printed slips kept literal: no printed word, symbol, hypothesis or number of the counted statement or of its proof was altered. Layout and class: the article's own class is \texttt{amsart} with packages including color, setspace, xcolor, here, graphicx, float, hyperref, tikz, fontenc, amsmath, amssymb, latexsym, amsthm, mathtools; the wrapper is the lane's pinned \texttt{amsart} editorial wrapper at 10pt on A4, which restates the theorem-like environments the retained text needs and adds the article's own macro definitions that the retained text needs (none), with extra wrapper packages (bm, bbm, dsfont, xspace, cancel, mathtools). The delivered numbering is the unit's own. Assets: no retained span contains a figure environment, a graphics-inclusion command, a PDF-page inclusion, a table or a TikZ picture; no image file is copied into this package, the package declares no asset dependency and no image file is redistributed with it. Licence: version arXiv:2603.21536v1 of this article is distributed under the Creative Commons Attribution 4.0 International licence, as recorded in the arXiv licence metadata for this exact version and shown on the abstract page https://arxiv.org/abs/2603.21536v1. A full rights notice is delivered at licenses/arxiv-2603.21536-CC-BY-4.0.txt. Characters: every retained excerpt is pure ASCII, so the delivered engine, XeLaTeX with its default Latin Modern text font, reports no missing character.
\begin{equation}\label{eq:gd-dR-fe-def-alt} 
\begin{cases} g_{0,0}(\varphi_{0}(x)) = \varphi_0 (f_{0,0}(x))  \\ g_{0,1}(\varphi_{1}(x)) = \varphi_0 (f_{0,1}(x))  \\ g_{1,0}(\varphi_{0}(x)) = \varphi_1 (f_{1,0}(x))  \\ g_{1,1}(\varphi_{1}(x)) = \varphi_1 (f_{1,1}(x))  \end{cases}   x \in [0,1]. 
\end{equation}

\begin{equation}\label{eq:def-f-dyadic} 
f_{0,0}(x) = f_{1,0}(x) = \frac{x}{2}, \ \ f_{0,1}(x) = f_{1,1}(x) = \frac{x+1}{2}, 
\end{equation}

\begin{equation}\label{eq:cd-posi} 
\min_{x \in [0,1]} c_{i,j} x + d_{i,j} = \min\{d_{i,j}, c_{i,j} + d_{i,j}\} \ge \sqrt{a_{i,j} d_{i,j} - b_{i,j} c_{i,j}} > 0, \ i, j \in \{0,1\}, 
\end{equation} 

\begin{equation}\label{eq:compatible-LF} 
0 = b_{i,0} < \dfrac{a_{i,0}}{c_{i,0} + d_{i,0}} = \dfrac{b_{i,1}}{d_{i,1}} < \frac{a_{i,1} + b_{i,1}}{c_{i,1} + d_{i,1}} = 1, \ \ i = 0,1.
\end{equation}  

\begin{equation}\label{eq:def-alpha} 
\alpha_i \coloneqq \frac{c_{i,0} + d_{i,0}}{a_{i,0}} - 2 = \frac{d_{i,1}}{b_{i,1}} - 2, \ i = 0,1. 
\end{equation} 

\begin{Thm}\label{thm:LF}
(i) If $\Phi \left(A_{i,j}^{\top}; \alpha_i \right) \ne \alpha_j$ holds for some $i, j \in \{0,1\}$, 
then, both of the solutions $\varphi_0$ and $\varphi_1$ of  \eqref{eq:gd-dR-fe-def-alt} are singular.\\
(ii) If $\Phi \left(A_{i,j}^{\top}; \alpha_i \right) = \alpha_j$ holds for every $i, j \in \{0,1\}$, 
then 
\begin{equation}\label{eq:solution-ac} 
\varphi_0 (x) = \frac{x}{-c_{0,0} x + 1 + c_{0,0}} \textup{ and } \varphi_1 (x) = \frac{x}{-c_{1,1} x + 1 + c_{1,1}}. 
\end{equation}
In particular, both $\varphi_0$ and $\varphi_1$ are smooth. 
\end{Thm}

\begin{proof}
(i) We show that $\varphi_0$ is singular. 

We first remark that  for every $y\in [0,1] \setminus D^{f}_0$ there exists a unique sequence $(i_n)_n$ in $\{0,1\}^{\mathbb N}$ such that $y \in I^{f}_{0} (i_1, \dots, i_n)$ for each $n \ge 1$. 

Assume that $x\in [0,1] \setminus D^{f}_0$ and $\varphi_0^{\prime}(x) > 0$.
Let $(i_n)_n$ be the unique sequence in $\{0,1\}^{\mathbb N}$ such that $x \in I^{f}_{0} (i_1, \dots, i_n)$ for each $n \ge 1$. 
Let 
\[ \begin{pmatrix} p_n & q_n \\  r_n & s_n \end{pmatrix} = A_{0,i_1} A_{i_1, i_2} \cdots A_{i_{n-1}, i_n}. \]
By the definition of $g_{i,j}$, it holds that 
\[ g_{0, i_1} \circ g_{i_1, i_2} \circ \cdots \circ g_{i_{n-1}, i_n} (x) = \Phi \left(A_{0,i_1} A_{i_1, i_2} \cdots A_{i_{n-1}, i_n}; x\right) = \frac{p_n x + q_n}{r_n x + s_n}, \ x \in [0,1]. \]

We show that $\min\{s_n, r_n + s_n\} > 0$ by induction on $n$. 
The case where $n=1$ is obvious. 
Assume that $\min\{s_n, r_n + s_n\} > 0$. 
Then $\displaystyle \min_{x \in [0,1]} r_n x + s_n = \min\{s_n, r_n + s_n\} > 0$. 
By \eqref{eq:cd-posi}, it holds that $d_{i_n, i_{n+1}} > 0$ and $c_{i_n, i_{n+1}} + d_{i_n, i_{n+1}} > 0$. 
By \eqref{eq:compatible-LF}, we see that 
$ \dfrac{b_{i_n, i_{n+1}}}{d_{i_n, i_{n+1}}} \in [0,1]$ and $\dfrac{a_{i_n, i_{n+1}} + b_{i_n, i_{n+1}}}{c_{i_n, i_{n+1}} + d_{i_n, i_{n+1}}}  \in [0,1]$. 
Since 
\[ (s_{n+1}, r_{n+1} + s_{n+1}) \]
\begin{equation}\label{eq:rs-iteration} 
 = \left(b_{i_n, i_{n+1}} r_{n} + d_{i_n, i_{n+1}} s_n,  (a_{i_n, i_{n+1}} + b_{i_n, i_{n+1}}) r_{n} + (c_{i_n, i_{n+1}} + d_{i_n, i_{n+1}}) s_n \right), 
\end{equation}
we obtain that 
\[ s_{n+1} = d_{i_n, i_{n+1}} \left(r_n \frac{b_{i_n, i_{n+1}}}{d_{i_n, i_{n+1}}} + s_n \right) > 0 \]
and 
\[ r_{n+1} + s_{n+1} = (c_{i_n, i_{n+1}} + d_{i_n, i_{n+1}}) \left(r_n \frac{a_{i_n, i_{n+1}} + b_{i_n, i_{n+1}}}{c_{i_n, i_{n+1}} + d_{i_n, i_{n+1}}} + s_n \right) > 0. \]

Since $b_{i,0} = 0 < d_{i,0}$, $ \Phi \left(A_{i, 0}^{\top} ; x\right)$ is well-defined for every $x \in \mathbb{R}$ and $i=0,1$. 
Since $\frac{r_n}{s_n} > -1 > -\frac{d_{i,1}}{b_{i,1}}$, $ \Phi \left(A_{i, 1}^{\top}; \frac{r_n}{s_n}\right)$ is well-defined for $i=0,1$ and $n \ge 1$. 
Therefore, we obtain that for every $n \ge 1$, 
\begin{equation}\label{eq:r-over-s} 
\frac{r_{n+1}}{s_{n+1}} = \Phi \left(A_{i_n, i_{n+1}}^{\top} ; \frac{r_n}{s_n} \right)
\end{equation}

By the definition of $f_{i,j}$ in \eqref{eq:def-f-dyadic}, it holds that $D^{f}_{0}$ is the set of dyadic rationals on $[0,1]$ and 
\[ \frac{|I^{g}_{0} (i_1, \dots, i_n)|}{|I^{f}_{0} (i_1, \dots, i_n)|} = 2^n \left( \frac{p_n + q_n}{r_n + s_n} - \frac{q_n}{s_n} \right) = 2^n \frac{p_n s_n - q_n r_n}{s_n (r_n + s_n)}. \]
Since 
$\displaystyle \lim_{n \to \infty} \frac{|I^{g}_{0} (i_1, \dots, i_n)|}{|I^{f}_{0} (i_1, \dots, i_n)|}  = \varphi_0^{\prime}(x) > 0$, 
we obtain that 
\[ \lim_{n \to \infty} \frac{|I^{g}_{0} (i_1, \dots, i_{n+1})| / |I^{f}_{0} (i_1, \dots, i_{n+1})| }{|I^{g}_{0} (i_1, \dots, i_n)| / |I^{f}_{0} (i_1, \dots, i_n)| } = 1. \]
This is equivalent to  
\[ \lim_{n \to \infty} \frac{\det(A_{i_n, i_{n+1}}) s_n (r_n + s_n)}{s_{n+1} (r_{n+1} + s_{n+1})} = \frac{1}{2}. \]

By \eqref{eq:rs-iteration}, this is equivalent to  
\[ \lim_{n \to \infty} t_n = \frac{1}{2}, \]
where we let 
\[ t_n \coloneqq \frac{\det(A_{i_n, i_{n+1}}) (\frac{r_n}{s_n} + 1)}{(b_{i_n, i_{n+1}} \frac{r_{n}}{s_n} + d_{i_n, i_{n+1}}) ((a_{i_n, i_{n+1}} + b_{i_n, i_{n+1}}) \frac{r_{n}}{s_n} + c_{i_n, i_{n+1}} + d_{i_n, i_{n+1}})}. \]
If $i_{n+1} = 0$, then $b_{i_n, i_{n+1}} = 0$ and hence 
\begin{equation}\label{eq:rs-zero} 
t_n = \frac{\frac{r_n}{s_n} + 1}{\frac{r_n}{s_n} + \frac{c_{i_n, 0} + d_{i_n, 0}}{a_{i_n, 0}}}. 
\end{equation} 
If $i_{n+1} = 1$, then $a_{i_n, i_{n+1}} + b_{i_n, i_{n+1}} = c_{i_n, i_{n+1}} + d_{i_n, i_{n+1}} > 0$ and hence 
\begin{equation}\label{eq:rs-one} 
t_n = \frac{\frac{d_{i_n,1}}{b_{i_n,1}} - 1}{\frac{r_n}{s_n} + \frac{d_{i_n, 1}}{b_{i_n, 1}}}. 
\end{equation}



For $y \in [0,1] \setminus  D_0^f$, let 
\[ \mathbb{N}_{i,j}(y) \coloneqq \left\{n \in \mathbb{N} \colon (i_n, i_{n+1}) = (i,j) \right\}, \ i, j \in \{0,1\},  \]
and 
\[ \mathbb{N}_{i}(y) \coloneqq \mathbb{N}_{i,0}(y) \cup \mathbb{N}_{i,1}(y), \ i \in \{0,1\}.  \]
Then $\{\mathbb{N}_{i,j}(y)\}_{i,j}$ are disjoint and $\mathbb{N} = \mathbb{N}_{0}(y) \cup \mathbb{N}_{1}(y)$. 
Let $\mathcal{N}$ be the set of $y \in [0,1] \setminus  D_0^f$ such that $\mathbb{N}_{i,j}(y)$ is finite for some $(i, j)$. 

Assume that $x \in [0,1] \setminus  \mathcal{N}$. 
Recall \eqref{eq:def-alpha}. 
Then, 
by \eqref{eq:rs-zero} and \eqref{eq:rs-one}, 
\[ \lim_{n \to \infty; n \in \mathbb{N}_i (x)} \frac{r_n}{s_n} = \alpha_i, \ i \in \{0,1\}.  \]
By this and \eqref{eq:r-over-s}, 
it holds that 
$\Phi \left(A_{i,j}^{\top}; \alpha_i \right) = \alpha_j$ for every $i, j \in \{0,1\}$. 
By the assumption, $x \in \mathcal{N}$. 
By the Borel normal number theorem, $\ell(\mathcal{N}) = 0$. 

Thus we see that $\varphi_0$ is singular. 
In the same manner, we also see that $\varphi_1$ is singular. 

(ii) We first remark that $\Phi(A;x) = \Phi(cA;x)$ for every $c \ne 0$ and every square matrix $A$. 
Since $a_{0,0}, a_{1,0}, b_{0,1}$ and $b_{1,1}$ are all positive, 
we can assume that $a_{0,0} = a_{1,0} = b_{0,1} = b_{1,1} = 1$. 

Since $\alpha_0 = c_{0,0} + d_{0,0} - 2$, 
the equation $\Phi \left(A_{0,0}^{\top}; \alpha_0 \right) = \alpha_0$ is equivalent to  the equation $(c_{0,0} + d_{0,0} -1)(d_{0,0} - 2) = 0$. 
Since $c_{0,0} + d_{0,0} > a_{0,0} = 1$,  we obtain that $d_{0,0} = 2$. 
By this and $b_{0,0} = 0$, 
we obtain that 
$A_{0,0} = \begin{pmatrix} 1 & 0 \\ c_{0,0} & 2 \end{pmatrix}$.  
Hence, 
$g_{0,0}(x) = \dfrac{x}{c_{0,0} x + 2}$. 


Since $\alpha_1 =  d_{1,1} - 2$, 
the equation $\Phi \left(A_{1,1}^{\top}; \alpha_1 \right) = \alpha_1$ is equivalent to  the equation $(c_{1,1} - d_{1,1} + 2)(d_{1,1} - 1) = 0$. 
Since $d_{1,1} > b_{1,1} = 1$, we obtain that $d_{1,1} = c_{1,1} + 2$. 
By this and $a_{1,1} + 1 = c_{1,1} + d_{1,1}$, 
we obtain that 
$A_{1,1} = \begin{pmatrix} 2c_{1,1} + 1 & 1 \\ c_{1,1} & c_{1,1} + 2 \end{pmatrix}$.  
Hence, 
$g_{1,1}(x) = \dfrac{(2c_{1,1} + 1)x + 1}{c_{1,1} x + c_{1,1} + 2}$. 

It also holds that $\alpha_i = c_{i,i}, \, i = 0,1$. 


We see that the equation $\Phi \left(A_{0,1}^{\top}; \alpha_0 \right) = \alpha_1$ is equivalent to  the equation $(c_{0,0} + 1)c_{0,1} + c_{0,0} (d_{0,1} - 1) = 2c_{1,1} (c_{0,0} + 1)$. 
By \eqref{eq:compatible-LF}, 
it holds that $d_{0,1} = c_{0,0} + 2$.  
Hence $c_{0,1} = 2 c_{1,1} - c_{0,0}$. 
Hence $a_{0,1} = c_{0,1} + d_{0,1} - 1 = 2c_{1,1} + 1$. 
Thus we obtain that $A_{0,1} = \begin{pmatrix} 2c_{1,1} + 1 & 1 \\ 2c_{1,1} - c_{0,0} & c_{0,0} + 2 \end{pmatrix}$.  
Hence, 
$g_{0,1}(x) = \dfrac{(2c_{1,1} + 1)x + 1}{(2c_{1,1}-c_{0,0}) x + c_{0,0} + 2}$. 

We see that the equation $\Phi \left(A_{1,0}^{\top}; \alpha_1 \right) = \alpha_0$ is equivalent to  the equation $c_{1,0} - c_{0,0} d_{1,0} = -c_{1,1}$. 
By \eqref{eq:compatible-LF}, 
it holds that $c_{1,0} + d_{1,0} = c_{1,1} + 2$.  
By solving the system of linear equations, we obtain that 
\[ c_{1,0} = \frac{c_{0,0} c_{1,1} + 2 c_{0,0} - c_{1,1}}{c_{0,0} + 1} \textup{ and }  d_{1,0} = \frac{2(c_{1,1} + 1)}{c_{0,0} + 1}. \]
By this and $b_{1,0} = 0$, 
we obtain that 
$A_{1,0} = \begin{pmatrix} 1 & 0 \\  \frac{c_{0,0} c_{1,1} + 2 c_{0,0} - c_{1,1}}{c_{0,0} + 1}  & \frac{2(c_{1,1} + 1)}{c_{0,0} + 1} \end{pmatrix}$.  
Hence, 
$g_{1,0}(x) = \dfrac{(c_{0,0} + 1)x}{(c_{0,0} c_{1,1} + 2 c_{0,0} - c_{1,1}) x + 2(c_{1,1} + 1)}$. 

Now by computations, we see that the functions $\varphi_0$ and $\varphi_1$ defined in \eqref{eq:solution-ac} satisfy \eqref{eq:gd-dR-fe-def-alt}.  
\end{proof}

\end{document}
